\chapter[Introdução]{Introdução}

Este capítulo apresenta os conceitos fundamentais que sustentam o desenvolvimento desta pesquisa, 
oferecendo ao(à) leitor(a) uma contextualização teórica acerca dos principais temas abordados ao longo do trabalho. 
Inicialmente, são discutidos aspectos relacionados à experimentação em Engenharia de Software e qualidade de produto.

A partir dessa fundamentação, são apresentados o delineamento do problema investigado e a questão de pesquisa que orienta este estudo.
Por fim, são descritos os procedimentos metodológicos adotados, incluindo a estratégia de pesquisa utilizada para condução da investigação, 
coleta e análise dos dados, de modo a garantir coerência científica e alinhamento aos objetivos propostos.

\section{Contexto}

\section{Problema}

\section{Questão de Pesquisa}

\section{Objetivos}

\section{Organização do Trabalho}

\section{Cronograma e Fluxo de Atividades}

\subsection{Atividades da Primeira Etapa}

\subsection{Atividades da Segunda Etapa}
